\section{Введение}

\subsection{Актуальность работы}

Выбор компаний для инвестирования предполагает оценку их перспектив по сведениям, доступным до получения результата. Исследование Gompers и соавторов, основанное на опросе венчурных инвесторов, показывает, что при отборе компаний существенное значение придаётся характеристикам команды и бизнеса~\cite{gompers2016}. Это придаёт прикладной смысл исследованию предикторов успеха: важно установить, какие наблюдаемые характеристики действительно помогают прогнозировать последующие результаты, а не только учитываются участниками рынка при принятии решений. Сам по себе опрос предпочтений инвесторов на этот вопрос не отвечает.

Различие между признаками, которым придают значение, и признаками, полезными для прогноза, определяет аналитическую проблему работы. Как показывает Shmueli, объяснение наблюдаемых зависимостей и прогнозирование новых наблюдений являются разными задачами: убедительность объяснительной модели не гарантирует её предсказательного качества~\cite{shmueli2010}. Поэтому содержательные предположения о роли команды, географии или направления деятельности нуждаются в самостоятельной прогностической проверке. Для исследования ключевых предикторов существенен не только факт связи признака с результатом, но и то, улучшает ли этот признак прогноз по сравнению с более простой моделью.

Открытые сведения позволяют поставить этот вопрос применительно к конкретному источнику данных. Каталог Y Combinator содержит карточки компаний со сведениями об их деятельности, командах и статусах~\cite{yc-directory}; на Kaggle опубликованы выгрузки этого каталога~\cite{kaggle-yc-dataset1,kaggle-yc-dataset2}. Доступность таких записей создаёт возможность эмпирической проверки, но сама по себе не подтверждает их пригодность для прогнозирования. С учётом разграничения описательного и прогностического анализа~\cite{shmueli2010} в данной работе проверяется, какую дополнительную информацию о последующем статусе компаний дают характеристики их публичных карточек сверх сведений о годе набора YC.

Таким образом, работа связывает два вопроса: какие характеристики компаний рассматриваются при оценке их перспектив~\cite{gompers2016} и какие из доступных характеристик выдерживают проверку на прогностическую полезность~\cite{shmueli2010}. Её предполагаемый вклад состоит в оценке дополнительной информативности сведений о команде, географии и деятельности компаний YC, а также в наглядном представлении надёжности полученных различий. Такой результат позволит судить о возможностях конкретного источника данных для прогнозирования наблюдаемого показателя успеха, а не предлагать универсальный перечень условий успешного предпринимательства.

\subsection{Цель}

Целью работы является выявление и визуализация ключевых предикторов наблюдаемого показателя успеха компаний Y Combinator на основе оценки их вклада в качество прогноза последующего состояния.

Успех в рамках исследования определяется как наличие статуса поглощённой или публичной компании в срезе от 1 августа 2026 года у компании, активной на начало наблюдения. Такое объединение выделяет два наблюдаемых корпоративных события, однако служит лишь приближённым показателем успеха: данные о статусе не позволяют оценить доходность для владельцев, а поглощение не обязательно означает выгодную продажу. Поэтому наряду с объединённым показателем рассматривается состав положительного класса. Продолжение деятельности и неактивность означают отсутствие данного показателя в позднем срезе, но не трактуются как равнозначные экономические результаты.

Предикторами выступают характеристики, зафиксированные на начало периода наблюдения, а ключевыми считаются те из них, которые устойчиво улучшают прогноз на данных, не использованных при обучении и настройке модели. Исследование оценивает прогностические связи, а не причинные эффекты; его выводы характеризуют полезность доступных сведений для прогноза, а не способы обеспечить успех компании.

\subsection{Объект исследования}

Объектом исследования являются компании каталога YC, имевшие активный статус в выгрузке от 13 июля 2023 года. Единицей наблюдения служит компания с уникальным идентификатором. Её характеристики на эту дату сопоставляются со статусом в версии позднего датасета от 1 августа 2026 года; срез от 27 февраля 2023 года используется для проверки согласованности сведений. Анализируются сторонние выгрузки каталога, размещённые на Kaggle, а не данные, собранные автором непосредственно с сайта YC.

Состав исследуемой совокупности определяется по июльскому срезу. Компании, отсутствующие в поздней выгрузке, сохраняются в учёте как записи с неизвестным последующим состоянием. Период сопоставления охватывает интервал с 13 июля 2023 года по 1 августа 2026 года, то есть немногим более трёх лет. Прогнозируется статус, зафиксированный в конечном срезе; точные даты переходов между состояниями по этим данным не определяются.

Выводы относятся к компаниям, активным на начало наблюдения. Измеренные характеристики отражают их состояние в июле 2023 года, а не обязательно на момент основания или поступления в акселератор. Поэтому результаты не переносятся автоматически на вновь созданные компании и предпринимательские инициативы вне YC.

\subsection{Предмет исследования}

Предметом исследования являются прогностические связи между характеристиками компаний на 13 июля 2023 года и их статусами на 1 августа 2026 года, а также вклад отдельных характеристик и их групп в качество прогноза наблюдаемого показателя успеха. Рассматриваются число основателей, размер команды, год основания, год и сезон набора YC, география, тематические теги и текстовые описания.

\subsection{Методы исследования}

\subsubsection{Методы предобработки данных}

Подготовка данных предусматривает проверку идентификаторов, устранение дубликатов, согласование форматов и сопоставление записей между срезами. Пропуски в признаках отделяются от неизвестных последующих состояний. Поздние значения не используются для восстановления характеристик на начало наблюдения; преобразования, параметры которых оцениваются по данным, настраиваются только на обучающей части выборки~\cite{kapoor2023}.

\subsubsection{Разведочный анализ данных (EDA)}

Разведочный анализ охватывает полноту данных, распределения характеристик, частоты статусов и связи между признаками. Отдельно сопоставляются компании разных годов набора и компании с известным и неизвестным последующим состоянием. Описательная статистика и графики служат для выявления особенностей выборки и ограничений дальнейшего моделирования.

\subsubsection{Статистический анализ и прогнозное моделирование}

Основной метод анализа --- сравнение моделей бинарной классификации наблюдаемого показателя успеха. Предусматривается сопоставление простых прогнозов, включая модель по году набора, с регуляризованной логистической регрессией и градиентным бустингом. Вклад групп признаков оценивается по изменению качества при их добавлении и исключении на одинаковых проверочных разбиениях.

Оценка качества включает способность различать классы и точность прогнозируемых вероятностей. Настройка моделей отделяется от итоговой проверки; различия в качестве сопровождаются оценками неопределённости. Основные проверки устойчивости касаются состава признаков, обработки пропусков и неизвестных последующих состояний. Отсутствие устойчивого улучшения относительно простого прогноза также является содержательным результатом.

\subsubsection{Анализ динамики}

Сопоставление срезов используется для описания переходов от активного состояния к наблюдаемым позднее статусам. Продолжающие работать компании учитываются наряду с поглощёнными, публичными и неактивными. Даты изменения статуса не приравниваются к датам выгрузок; по двум наблюдениям не восстанавливается полная история компании.

\subsubsection{Инструменты визуализации}

Для визуализации предусмотрены библиотеки Python Matplotlib и Seaborn. Графики представляют состав выборки, переходы между статусами, сравнение моделей и вклад признаков с неопределённостью оценок. Подписи указывают анализируемую совокупность, единицы измерения и преобразования, необходимые для самостоятельного прочтения изображения.

\subsection{Задачи}

\begin{enumerate}
    \item Рассмотреть исследования предикторов успеха стартапов и подходы к определению прогнозируемого результата.
    \item Описать происхождение и временные ограничения данных, сформировать выборку для сопоставления характеристик и последующих статусов.
    \item Провести разведочный анализ, оценить полноту признаков и особенности записей с неизвестным последующим состоянием.
    \item Сравнить простые и многомерные прогнозные модели, оценить дополнительную информативность сведений о команде, географии и деятельности компаний.
    \item Проверить устойчивость различий в качестве и представить вклад признаков вместе с неопределённостью оценок.
    \item Визуализировать результаты и сформулировать выводы о возможностях прогнозирования по данным каталога и границах применимости этих выводов.
\end{enumerate}

\subsection{План работы}

В таблице~\ref{tab:plan} приведён первоначальный календарный план с сохранением исходных задач и сроков. Он отражает прежнюю постановку исследования: упоминания гипотез H1--H8 и периода 2023--2025 годов относятся к первоначальному замыслу, а не к обновлённому анализу.

\begin{table}[htbp]
\centering
\small
\caption{Календарный план проекта (25.01.2026\,--\,01.04.2026)}
\label{tab:plan}
\begin{tabular}{@{}>{\raggedright\arraybackslash}p{3.2cm} >{\raggedright\arraybackslash}p{2.3cm} >{\raggedright\arraybackslash}p{6.8cm} >{\raggedright\arraybackslash}p{3.5cm}@{}}
\toprule
\textbf{Этап} & \textbf{Неделя / Период} & \textbf{Основные задачи} & \textbf{Результат} \\
\midrule

\textbf{Этап 1: Подготовка данных и EDA} & Неделя 1 \newline 25.01--31.01 & Сбор данных (Kaggle); настройка окружения; слияние источников; обработка пропусков & Объединенные датасеты \\
\addlinespace[3pt]
-- & Неделя 2 \newline 01.02--07.02 & Создание производных признаков (бинарный статус, тип основателей); EDA: распределения, пропуски & Датасет A и B + ноутбук \\
\addlinespace[3pt]
-- & Неделя 3 \newline 08.02--14.02 & Анализ категориальных признаков; визуализация батчей, отраслей, стран & EDA-отчет с графиками \\
\midrule

\textbf{Этап 2: Анализ предикторов} & Неделя 4 \newline 15.02--21.02 & Проверка гипотез H1--H4: соло vs команда, B2B vs Consumer, сезон батча, лаконичность описания & Графики и статистические тесты \\
\addlinespace[3pt]
-- & Неделя 5 \newline 22.02--28.02 & Проверка гипотез H5--H8: количество тегов, найм, география, стадия развития & Полный набор визуализаций \\
\addlinespace[3pt]
-- & Неделя 6 \newline 01.03--07.03 & Корреляционный анализ; анализ динамики 2023--2025; матрица переходов & Корреляционная матрица + выводы \\
\midrule

\textbf{Этап 3: Оформление} & Неделя 7 \newline 08.03--14.03 & Написание отчета: введение, описание данных, методология & Черновик курсовой работы \\
\addlinespace[3pt]
-- & Неделя 8 \newline 15.03--21.03 & Написание отчета: результаты, выводы, оформление графиков & Полный текст работы \\
\addlinespace[3pt]
-- & Финал \newline 22.03--01.04 & Проверка оформления; создание презентации; финализация & Готовый проект к защите \\
\bottomrule
\end{tabular}
\end{table}

Таблица~\ref{tab:revision-plan} содержит план переработки с 12 августа 2026 года до защиты 3 октября 2026 года. Срок сдачи отчётных документов, включая текст курсовой работы, --- 18 сентября; презентация должна быть подготовлена и сдана до 1 октября включительно. Редактура текста ведётся параллельно с анализом, чтобы завершить проверку результатов до сдачи документов. План задаёт сроки задач, а не подтверждает их выполнение.

\begin{table}[htbp]
\centering
\small
\caption{План переработки и подготовки к защите (12.08.2026\,--\,03.10.2026)}
\label{tab:revision-plan}
\begin{tabular}{@{}>{\raggedright\arraybackslash}p{3.2cm} >{\raggedright\arraybackslash}p{2.3cm} >{\raggedright\arraybackslash}p{6.8cm} >{\raggedright\arraybackslash}p{3.5cm}@{}}
\toprule
\textbf{Этап} & \textbf{Период} & \textbf{Основные задачи} & \textbf{Результат} \\
\midrule
\textbf{1. Постановка исследования} & 12.08--18.08 & Проверка литературы; уточнение показателя успеха, совокупности и периода наблюдения; редактура введения & Постановка и проверенные источники \\
\addlinespace[3pt]
\textbf{2. Подготовка данных} & 19.08--25.08 & Сопоставление компаний; разделение исходных признаков и поздних статусов; учёт пропусков и ненайденных записей & Аналитическая выборка и отчёт о преобразованиях \\
\addlinespace[3pt]
\textbf{3. Обновление EDA} & 26.08--01.09 & Пересчёт распределений и переходов; анализ пропусков и когорт; обновление графиков & Графики и раздел EDA \\
\addlinespace[3pt]
\textbf{4. Прогнозные модели} & 02.09--08.09 & Фиксация схемы проверки и метрик; сравнение простых прогнозов, логистической регрессии и градиентного бустинга & Сравнение моделей без утечки данных \\
\addlinespace[3pt]
\textbf{5. Предикторы и устойчивость} & 09.09--14.09 & Оценка вклада групп признаков, неопределённости и качества вероятностей; проверка чувствительности к пропускам & Оценки информативности и устойчивости \\
\addlinespace[3pt]
\textbf{6. Переработка текста} & 12.08--16.09 \newline Параллельно & Обновление обзора, методов и результатов; согласование чисел и графиков; исправление неподтверждённых выводов & Полный текст курсовой работы \\
\addlinespace[3pt]
\textbf{7. Итоговая проверка} & 15.09--17.09 & Повторный запуск анализа; проверка ссылок и сборки LaTeX; подготовка комплекта отчётных документов & Проверенный комплект документов \\
\addlinespace[3pt]
\textbf{8. Сдача документов} & \textbf{18.09} & Сдача отчётных документов, включая текст курсовой работы & Документы сданы \\
\addlinespace[3pt]
\textbf{9. Презентация} & 19.09--01.10 & Подготовка слайдов и доклада по сданной работе; проверка соответствия результатов; сдача презентации до 01.10 включительно & Презентация сдана \\
\addlinespace[3pt]
\textbf{10. Подготовка к выступлению} & 02.10 & Репетиция доклада; подготовка ответов на вопросы о методах и ограничениях & Готовность к выступлению \\
\addlinespace[3pt]
\textbf{11. Защита} & \textbf{03.10} & Доклад и ответы на вопросы комиссии & Защита курсовой работы \\
\bottomrule
\end{tabular}
\end{table}
\clearpage
